\documentclass[12pt, a4paper]{article} % Classe: article, report o book [8]

% --- Lingua e Codifica (Supporto Italiano) ---
\usepackage[utf8]{inputenc}  % Codifica input
\usepackage[T1]{fontenc}      % Codifica font
\usepackage[italian]{babel}   % Lingua italiana [2]

% --- Impaginazione e Margini ---
\usepackage[margin=2.5cm]{geometry} % Imposta margini (2.5cm su tutti i lati)

% --- Pacchetti Scientifici/Matematici ---
\usepackage{amsmath, amssymb, amsfonts, amsthm} % Matematica avanzata [2]
\usepackage{graphicx} % Per inserire immagini
\usepackage{hyperref} % Per collegamenti ipertestuali
\usepackage{booktabs} % Per tabelle professionali

% --- Formattazione Testo e Altro ---
\usepackage[document]{ragged2e}
\usepackage{xcolor}    % Per usare i colori 

% Document metadata
\title{Interviste}
\author{Mattia Cavina}
\date{10 Dicembre 2025}

\begin{document} 

\maketitle
\begin{abstract} 
      In questo documento sono riportate le integrazioni alle interviste
      precedentemente effettuate con i repparti non dierttamnete coinvolti nell'analisi,
      ma utili a completare le informazioni raccolte.
\end{abstract}
\section{Ufficio Project Management}
\paragraph{Ruolo}
L'ufficio project management si occupa della pianificazione e gestione delle
commesse, coordinando le attività tra il reparto commerciale e i reparti
tecnici.
\paragraph{Richieste}
Le richieste principali dell'ufficio project management riguardano:
\begin{itemize}
      \item \textbf{Gestione contratto tecnico:} Si sono resi disponibili a stilare un contratto tecnico in parallelo a quello commerciale, che contenga tutte le variazioni contrattuali attuate per il corretto funzionamnento del progetto.
      \item \textbf{Costificazione variazione:} Hanno espresso la possibilità di costificare le variazioni tecniche in modo separato dal contratto commerciale, per valutare il margine di manovra sulle modifiche necessarie
\end{itemize}

\section{Ufficio Manualistica}
\paragraph{Ruolo}
L'ufficio manualistica si occupa della creazione e gestione dei manuali tecnici
e delle certificazioni dei prodotti venduti.
\paragraph{Richieste}
Le richieste principali dell'ufficio manualistica riguardano:
\begin{itemize}
      \item \textbf{Campo Certificazione:} Prevedere un campo nel contratto dove indicare le certificazioni richieste dal componente venduto (CE, ATEX, FDA, ecc \dots).
      \item \textbf{Prodotti di Acquisto:} Prevedere un campo Norme di riferimento per i prodotti acquistati da terzi che hanno delle certificazioni specifica (ATEX, ALIMENTARI, ecc \dots).
      \item \textbf{Caratteristiche Tecniche:} Dare importanza alle caratteristiche tecniche dei componenti venduti per facilitare la creazione dei manuali ( dati elettrici per un motore, dati meccanici per macchinario, allestimento, accessori previsti, ecc \dots).
\end{itemize}

\section{Ufficio Acquisti}
\paragraph{Ruolo}
L'ufficio acquisti si occupa della gestione degli ordini di acquisto per i
componenti necessari alla produzione. Tra i compiti vi è anche il controllo
delle richieste di assistenza in funzione dell'acquisto dei ricambi. In questa
fase consultano i contratti generati dai commerciali per verificare le parti da
ordinare ed eventuali discripanze tecniche o commerciali.

\end{document}